% =====================================================================
%  04_metodologia.tex — Prototipado evolutivo y plan de prototipos
% =====================================================================
\chapter{Metodología}
\label{cap:metodologia}

\section{Prototipado evolutivo}

El desarrollo de \saru{} sigue la metodología de \textbf{prototipado
evolutivo}: el sistema se construye mediante versiones funcionales
sucesivas, cada una de las cuales se entrega, se prueba y se usa como base
de la siguiente~\cite{sommerville2016}. A diferencia del prototipo
desechable, cada versión no se descarta: se corrige y se amplía.

Esta metodología es adecuada para \saru{} por tres razones:

\begin{itemize}
  \item \textbf{Los requerimientos cambian con el uso.} Muchos casos solo
        aparecen con video real del patio: unidades detenidas sobre la
        línea, varias unidades a la vez o números leídos con dígitos de
        menos. Cada prototipo en sitio revela casos nuevos.
  \item \textbf{Los componentes son heterogéneos.} Cámara, modelos de
        visión, lógica de cruce, base de datos, mensajería y web dependen
        unos de otros; validarlos juntos en una sola entrega es riesgoso.
  \item \textbf{El sistema ya opera.} Cada prototipo debe mantener el
        servicio funcionando mientras se mejora.
\end{itemize}

\section{Ciclo de cada prototipo}

Cada prototipo recorre las mismas etapas y se documenta con la misma
estructura en el capítulo~\ref{cap:prototipos}:

\begin{enumerate}
  \item \textbf{Especificación}: requerimientos que atiende e historias de
        usuario con sus criterios de aceptación~\cite{cohn2004}.
  \item \textbf{Diseño}: arquitectura, componentes e interfaz.
  \item \textbf{Desarrollo}: implementación y despliegue en la Jetson.
  \item \textbf{Pruebas}: protocolo con casos identificados por módulo y
        evidencia de cada resultado.
  \item \textbf{Retroalimentación}: los fallos y los casos nuevos se
        convierten en requerimientos del siguiente prototipo.
\end{enumerate}

\section{Plan de prototipos}

\begin{figure}[htbp]
  \centering
  \resizebox{\linewidth}{!}{%
  \begin{tikzpicture}[node distance=0.6cm]
    \node[caja, minimum width=3.2cm] (p1) {\textbf{Prototipo 1}\\Sistema base\\en sitio};
    \node[cajaGris, minimum width=3.2cm, right=of p1] (p2) {\textbf{Prototipo 2}\\Operación\\continua};
    \node[cajaGris, minimum width=3.2cm, right=of p2] (p3) {\textbf{Prototipo 3}\\Mejora de\\la lectura};
    \node[cajaAlt, minimum width=3.2cm, right=of p3] (pf) {\textbf{Prototipo final}\\Integración\\y validación};
    \draw[flecha] (p1) -- (p2);
    \draw[flecha] (p2) -- (p3);
    \draw[flecha] (p3) -- (pf);
    \node[etiqueta, below=0.25cm of p1] {Reportado en este documento};
    \node[etiqueta, below=0.25cm of p2] {Propuesta};
    \node[etiqueta, below=0.25cm of p3] {Propuesta};
    \node[etiqueta, below=0.25cm of pf] {Antes del cierre};
  \end{tikzpicture}%
  }
  \caption{Prototipado general del proyecto.}
  \label{fig:prototipado}
\end{figure}

\begin{description}[style=nextline]
  \item[Prototipo 1: sistema base en sitio] La cadena completa de
        detección, lectura, cruce, registro, avisos, tablero web y
        grabación continua, operando en vivo en el patio. Es el que reporta
        este documento (sección~\ref{sec:prototipo1}).
  \item[Prototipo 2: operación continua] Llevar el sistema a operación
        permanente en el gabinete del patio con una Jetson de
        16~GB~\cite{rep-tecnico}: arranque automático, reinicio ante fallas
        y verificación de que los avisos se entregan.
  \item[Prototipo 3: mejora de la lectura] Ampliar los datos de
        entrenamiento con los recortes del patio, reentrenar los modelos
        donde la lectura falla (por ejemplo, vistas no laterales) y medir
        la exactitud de cada etapa.
  \item[Prototipo final: integración y validación] Integrar los cambios de
        los prototipos anteriores y validar el sistema de extremo a
        extremo con pruebas en tiempo real, antes del cierre del proyecto.
\end{description}

\begin{aviso}[title={Por validar}]
El contenido de los prototipos 2 y 3 es una propuesta basada en los
pendientes registrados en los reportes del equipo. Debe confirmarse con la
coordinación del proyecto.
\end{aviso}
